\documentstyle[% 


righttag% 


,11pt% 


,amssymb% 


,russian%               remove in Eng. 


,izvru%                 Set izven% in Eng. 


% ,izvru-ks             for brief communications 


]{amsart} 


 


%       Math definitions 


 


\newcommand{\Tr}{\operatorname{Tr}}


\newcommand{\Ker}{\operatorname{Ker}}


\newcommand{\ad}{\operatorname{ad}}


\newcommand{\pr}{\operatorname{pr}}


\newcommand{\ort}{\operatorname{ort}}


\newcommand{\gr}{\operatorname{gr}}


\newcommand{\id}{\operatorname{id}}


\newcommand{\ol}[1]{\overline {#1}}


\newcommand{\olh}[1]{\overline {#1}{}^h}


\newcommand{\olv}[1]{\overline {#1}{}^v}


\newcommand{\g}[1]{\widetilde {#1}}


\newcommand{\grh}[1]{\widetilde {#1}{}^h}


\newcommand{\grv}[1]{\widetilde {#1}{}^v}


\newcommand{\gh}[1]{{#1}{}^h}


\newcommand{\gv}[1]{{#1}{}^v}


\newcommand{\1}{\frac12}


\newcommand{\4}{\frac14}


\newcommand{\8}{\frac18}


\newcommand{\6}{\frac1{16}}


\newcommand{\chp}[1]{\frac\partial{\partial #1}}


\newcommand{\ch}[1]{\frac\partial{\partial #1^t}}


\newcommand{\w}{\omega}


\newcommand{\n}[2]{\nabla _{#1}{#2}}


\newcommand{\grn}[2]{\widetilde \nabla _{#1}{#2}}


\newcommand{\tri}{\triangle}


\newcommand{\thX}{\tilde X{}^h}


\newcommand{\tvQ}{\tilde Q{}^v}


\newcommand{\tvS}{\tilde S{}^v}


\newcommand{\thY}{\tilde Y{}^h}


\newcommand{\tvT}{\tilde T{}^v}


\newcommand{\thF}{\tilde F{}^h}


\newcommand{\tts}[1]{}


 


\theoremstyle{plain} 


\theorembodyfont{\it} 


\newtheorem{props}{\hskip\parindent Предложение}[section]


\newtheorem{lems}{\hskip\parindent Лемма}[section]


 


\theoremstyle{definition} 


\theorembodyfont{\rm} 


\newtheorem{cors}{\hskip\parindent Следствие}[section]


\newtheorem{defns}{\hskip\parindent Определение}[section]


 


%\hoffset=-15mm                  For Eng. 


 


\setcounter{page}{57} 


\god{1997} 


\nomer{9~(424)} %                Remove in Eng. 


%\nomer{Vol.\,41, No.\,9}        Set in Eng. 


% \str{75-81}                    Set in Eng. 


\udk{514.76} 


\author{\it Н.К.\,Фарафонова} 


\title{Геометрия грассманова расслоения. I} 


 


\date{} 


% \pagestyle{myheadings}         Set in Eng. 


 


\pagestyle{plain} %              Remove in Eng. 


\begin{document} 


 


% \thispagestyle{plain}          Set in Eng. 


 


\maketitle 


 


% Body text 


 





\section*{Введение}


Изучение геометрии тотального пространства расслоения приобрело самостоятельное


значение после работ 50-х годов Б.Л.\,Лаптева и Сасаки, история этих вопросов


освещена в обзорах \cite{Shap} и \cite{Bor}.


Целью этой статьи является установление связи между геометрическими


характеристиками грассманова расслоения и римановым многообразием служащим его


базой.


При этом грассманово расслоение наделяется естественной римановой


метрикой, геометрический смысл которой был исследован автором в \cite{Far2}.


Аналогичная задача для касательного, сферического и нормального расслоений,


а также расслоения реперов решалась Сасаки, Домбровским, Моком, Ямпольским


и другими авторами. Подробная информация о полученных по этой проблеме


результатах имеется в третьем параграфе \cite{Bor}.


В первой части статьи строится образующая система векторных полей грассманова


расслоения, свойства которой позволяют связать геометрические характеристики


базы и тотального пространства расслоения.


Основной технической трудностью, которую пришлось преодолеть для грассманова


расслоения, была задача построения специальных векторных полей грассманова


расслоения, являющихся лифтами кососимметрических тензорных полей (определение


2.2). Всюду в дальнейшем римановы многообразия предполагаются гладкими, класса


$C^{\infty}$, как само многообразие, так и его риманова метрика, и всюду


действует правило Эйнштейна.





\section{Кососимметрические тензоры на римановом многообразии}





Этот параграф посвящен краткому изложению как хорошо известных, так


и специально здесь вводимых понятий и свойств для (1,1)-тензоров на римановом


многообразии $М$ с римановым метрическим тензором $g$, которые будут широко


использоваться в качестве технических средств в \S\,3 для описания


геометрии грассманова расслоения. Наиболее подходящей нашим целям


интерпретацией 1-контравариантного и 1-ковариантного тензорного поля будет


служить операторная интерпретация (\cite{Kob1}, гл.\,1, пр.\,3.1),


согласно которой


тензорное поле $Q$ типа (1,1) является линейным отображением модуля гладких


векторных полей на $М$ в себя. При этом (1,1)-тензор в точке $u_0=(u_0^a)$


многообразия $М$ является линейным оператором


\begin{equation}


Q=Q^b_adu^a\otimes \frac{\partial}{\partial u^b}:T_{u_0}M\to T_{u_0}M


\end{equation}


в касательном пространстве. Компоненты этого тензора $Q^b_a(u)$ относительно


координат $u^b$ являются гладкими функциями этих координат. (1,1)-тензорное


поле $Q$ называется кососимметрическим на многообразии $M$, если оно


кососимметрично в каждом касательном пространстве относительно скалярного


произведения, индуцированного римановой метрикой


\begin{equation}


g(QX,Y)+g(X,QY)=0,


\end{equation}


где $X$ и $Y$ --- любые касательные векторы к $M$.


В дальнейшем, если не оговорено противное, будем называть кососимметрические


тензорные поля типа (1,1) просто кососимметрическими, и если из контекста


понятно, относительно какого базиса рассматривается матрица линейного


оператора или квадратичная форма, то их матрица будет обозначаться той же


буквой, что и оператор или квадратичная форма.





Мы будем рассматривать компоненты $S^b_a$ (1,1)-тензора $S$ и относительно


произвольного (не обязательно координатного) поля реперов $E=(E_1,\dotsc,E_n)$,


однозначно определяемые формулой


\begin{equation}


S=S^b_a\chi ^a\otimes E_b,


\end{equation}


где $\chi=(\chi^a)$ --- кореперное поле двойственное к $E$, аналогичной формуле


(1).





Классическое отождествление $\boldsymbol\Lambda^2(M)$


с пространством $E(M)$ кососимметрических эндоморфизмов для $T_uM$,


(\cite{Kob1}, c.\,152; \cite{Kob2}, с.\,317)


позволяет для любой пары векторных полей $X$ и $Y$


на многообразии $M$ определить формулой


\begin{equation}


X\land Y(U):=\frac{1}{2}(g(U,Y)X-g(U,X)Y)


\end{equation}


кососимметрическое тензорное поле, отвечающее разложимому бивектору $X\land Y$.


Очевидно, что в некоторой окрестности каждой точки многообразия


существует базис в модуле кососимметрических тензорных полей, составленный из


$\frac{n(n-1)}{2}$ кососимметрических тензорных полей $U_a\land U_b$,


отвечающих разложимым бивекторам. Например, выберем в качестве $U_a$


координатные векторы $\frac{\partial}{\partial u_a}$. В дальнейшем для


краткости будем называть кососимметрические тензорные поля, отвечающие


разложимому бивектору, просто разложимыми бивекторами.


В дальнейшем скобочная операция от кососимметрических тензоров --- это всегда


операторный коммутатор, а от векторных полей --- это скобка Ли.


Следующее предложение описывает операторные свойства разложимых бивекторов.





\begin{props} Пусть $Q$ --- произвольное кососимметрическое


тензорное поле, а $X$, $Y$, $Z$ --- произвольные векторные поля на римановом


многообразии $M$. Тогда


\begin{itemize}


\item[а)] $[Q,X\land Y]=QX\land Y+X\land {QY};$


\item[б)] $\underset{(X,Y,Z)}{S}{}X\land Y(Z)=X\land Y(Z)+


Y\land Z(X)+Z\land X(Y)=0.$


\end{itemize}


\end{props}





\begin{cors} Если векторные поля $X$, $Y$, $U$, $W$ на римановом


многообразии $M$ таковы, что $X$ и $Y$ ортогональны $U$ и $W$, то





a) $[X\land W,Y\land W]=-\frac12\| W\| ^2X\land Y$,


\qquad


б) $[X\land Y,U\land W]=0$.


\end{cors}





Пользуясь стандартным определением продолжения риманова метрического тензора


$g$ на всю алгебру тензоров $\Bbb T (M)$ (\cite{Kob1}, c.\,152),


приведем метрические свойства кососимметрических тензоров.





\begin{props} Для любых кососимметрических тензорных полей $P$,


$Q$ и $S$ и любых векторных полей $X$, $Y$ на римановом многообразии $M$


имеют место следующие равенства: {\rm{a)}} $g(Q,S)=-\Tr QS$,


{\rm{б)}} $g(Q,X\land Y)=g(QY,X)$, {\rm{в)}} $g([Q,S],P)+g(S,[Q,P])=0$,


где $\Tr$ --- следный функционал.


\end{props}





{\bf Доказательство} a) (\cite{Kob1}, гл.\,4, \S\,1, пример 1.6).


б) В плоскости,


натянутой на векторы $X$ и $Y$, выберем орты $U$ и $V$.


Дополним их ортонормальными векторами до базиса в касательном пространстве.


Причем на этих дополнительных векторах разложимый бивектор $X\land Y$


действует тривиально. Тогда, воспользовавшись равенством а) и тем, что след


не зависит от базиса, в котором он вычисляется, получим


\begin{gather*}


g(Q,X\land Y)=-\Tr Q\circ X\land Y= -(g(Q\circ X\land Y(U),U)+g(Q\circ X


\land Y(V),V))=\\


=-\frac12 (g(U,Y)g(QX,U)-g(U,X)g(QY,U)+g(V,Y)g(QX,V)-\\


-g(V,X)g(QY,V))=-\frac12 (g(g(QX,U)U+g(QX,V)V,Y)-g(g(QY,U)U+\\


+g(QY,V)V,X))=-\frac12(g(QX,Y)-g(QY,X))=g(QY,X).


\end{gather*}


в) Из свойства а) и инвариантности следного функционала


имеем


$$


g([Q,S],P)+g(S,[Q,P])=-\Tr [Q,S]P-\Tr S[Q,P]=0.\quad\square


$$





\begin{cors} $g(X\land Y,X\land Y)=\frac12 (\| X\| ^2


\| Y\| ^2-g(X,Y)^2)$.


\end{cors}





Пусть $\nabla_X$ --- ковариантная производная римановой связности на $M$,


которая стандартным образом определяется на модуле всех тензорных полей и


сохраняет кососимметричность (1,1)-тензоров. В дальнейшем


будет использоваться матричное равенство для ковариантной производной


(1,1)-тензорного поля $Q$ вдоль векторного поля $Y$, заданных матрицей и


вектор-столбцом их компонент относительно произвольного реперного поля $E$


\begin{equation}


\nabla_YQ=dQ(Y)+\Gamma (Y)Q-Q\Gamma (Y),


\end{equation}


где $\Gamma (Y)=\Gamma (\chi )(Y)$ --- матрица значений компонент матрицы


Кристоффеля (\cite{Far1}, \S\,2) относительно реперного поля $E$, т.\,е.


значений


1-форм, на векторе $Y$. Дифференциал матрицы $dQ$ понимается, как обычно,


покомпонентно, т.\,е. это матрица 1-форм на многообразии $M$. Для


матрицы Кристоффеля $\Gamma(\chi)$ выполняется следующее часто


используемое свойство: относительно ортонормального реперного


поля $E$ она является кососимметрической, т.\,е.


$\Gamma^t(\chi)=-\Gamma(\chi)$.





Пусть $R(X,Y)$ --- преобразование кривизны, которое мы будем рассматривать


продолженным на всю алгебру тензоров $\Bbb T(M)$ на римановом


многообразии $M$:


\begin{equation}


R(X,Y)Q=\nabla_X\nabla_YQ-\nabla_Y\nabla_XQ-\nabla_{[X,Y]}Q.


\end{equation}


С учетом отождествления (4) пространства бивекторов с пространством


кососимметрических тензоров определим оператор кривизны $\rho$, как


самосопряженный эндоморфизм в пространстве кососимметрических


тензоров формулой $\rho(X\land Y):=R(X,Y)$. Самосопряженность, выраженная


формулой


\begin{equation}


g(R(X,Y)Z,W)=g(R(Z,W)X,Y),


\end{equation}


следует из


последнего в ряду четырех алгебраических свойств $R(X,Y)$, приведенных в


(\cite{Kob1}, гл.\,V). Очевидно, что на основании


предложения 1.2~б) для оператора кривизны выполняется свойство


\begin{equation}


g(\rho(X\land Y),U\land V)=g(R(X,Y)V,U).


\end{equation}





\begin{props} Для любых кососимметрических тензоров $Q$ и $S$ на римановом


многообразии


$M$ и любых векторов $X$, $Y$ и $Z$ выполняются равенства:


\begin{alignat*}{2}


&a)\ g(\rho(Q),S)=g(Q,\rho(S)), &\qquad &б)\ g((\nabla_Z\rho)(Q),S)=


g(Q,(\nabla_Z\rho)(S)),\\


&в) \ R(X,Y)Q=[R(X,Y),Q], &\qquad &г) \ \nabla_X[Q,S]=[\nabla_XQ,S]+[Q,\nabla_XS].


\end{alignat*}


\end{props}





{\bf Доказательство.}


В силу линейности отображений рассматриваемых в свойствах а)--в), достаточно


доказать эти равенства для случая, когда кососимметрические тензоры $Q$ и $S$


отождествляются по формуле (4) с разложимыми бивекторами, которые образуют


базис.





а) Из формул (7) и (8) следует, что $g(\rho(X\land Y),U\land V)=g(R(X,Y)V,U)=


g(R(V,U)X,Y)=g(X\land Y,\rho(U\land V))$.





б) Из определения ковариантной производной от тензорного поля


\begin{gather*}


g((\nabla_Z\rho)(X\land Y),U\land V)=g(\nabla_Z\rho(X\land Y),U\land V)-


g(\rho(\nabla_Z(X\land Y)),U\land V)=\\


=g(\nabla_Z\rho(X\land Y),U\land V)-g(\rho(\nabla_ZX\land Y),U\land V)-


g(\rho(X\land {\nabla_ZY}),U\land V)=\\


=g(\nabla_ZR(X,Y),U\land V)-g(R(\nabla_ZX,Y)V,U)-g(R(X,\nabla_ZY)V,U)=


\end{gather*}


Так как $R(X,Y)$ --- кососимметрический тензор, а ковариантная производная


сохраняет кососимметричность тензоров (1,1) и наследует все алгебраические


свойства тензора кривизны, то воспользовавшись предложением 1.2~б) и формулой


(7), продолжим равенство


\begin{gather*}


=g(\nabla_ZR(X,Y)V-R(\nabla_ZX,Y)V-R(X,\nabla_ZY)V,U)=g((\nabla_ZR)


(X,Y)V,U)=\\


=g((\nabla_ZR)(V,U)X,Y)=g((\nabla_Z\rho)(V\land U),Y\land X)=


g((\nabla_Z\rho)(U\land V),X\land Y).


\end{gather*}


в) Необходимо отметить, что слева в равенстве п.\,в) стоит оператор


преобразования кривизны из (6), а справа $R(X,Y)$ рассматривается, как


кососимметрический тензор. Используя предложение 1.1~а) и формулу


\begin{equation}


\nabla_X(U\land V)=\nabla_XU\land V+U\land {\nabla_XV},


\end{equation}


которая легко может быть получена из формулы для ковариантной производной


от тензорного поля (\cite{Kob1}, гл.\,3, предл.\,2.10), получим равенство


$$


R(X,Y)(U\land V)=[R(X,Y),U\land V].


$$


г) Из определения коммутаторного оператора и


(\cite{Kob1}, гл.\,3, предл.\,2.10)


\begin{gather*}


(\nabla_X[Q,S])(U)=\nabla_X([Q,S]U)-[Q,S]\nabla_XU=\nabla_X(QS(U)-SQ(U))-\\


-QS(\nabla_XU)+SQ(\nabla_XU)=\nabla_XQS(U)+Q\nabla_XS(U)-\nabla_XSQ(U)-


S\nabla_XQ(U)-\\


-QS(\nabla_XU)+SQ(\nabla_XU)=[\nabla_XQ,S](U)+[Q,\nabla_XS](U).\quad\square


\end{gather*}


\medskip





Пусть $\Pi$ --- некоторое $l$-мерное подпространство в касательном пространстве


$T_wM$, $w\in M$. Пусть (1,1)-тензоры $\pr$, $\ort $ --- это операторы в $T_wM$


проектирования на $\Pi$ и на ортонормальное дополнение к $\Pi$,


соответственно.





\begin{defns}


(1,1)-тензоры, определяемые формулами


\begin{align}


|Q|_\Pi &=\pr Q \pr +\ort Q\ort,\\


\{ Q\} _\Pi &=\pr Q\ort +\ort Q\pr ,


\end{align}


называются внутренней и внешней компонентами (1,1)-тензора $Q$


соответственно относительно подпространства $\Pi$.


\end{defns}





Когда из контекста понятно, относительно какого подпространства


$\Pi$ берутся внешняя или внутренняя компоненты тензора, то индекс $\Pi$


часто опускается.





\begin{defns}


Подпространства кососимметрических тензоров, имеющих нулевую внешнюю


относительно


$\Pi$ компоненту, обозначаются $I(\Pi)$, а нулевую внутреннюю --- $E(\Pi)$ и


называются соответственно внутренним и внешним относительно $\Pi$


подпространством кососимметрических тензоров.


\end{defns}





\begin{props}


Всякий $(1,1)$-тензор $Q$ в точке $u$ многообразия $M$ единственным образом


разлагается по формуле


\begin{equation}


Q=|Q|_\Pi +\{ Q\} _\Pi,


\end{equation}


где $\Pi$ --- произвольное подпространство в $T_wM$. Если $Q$ ---


кососимметрический тензор, то его внутренняя и внешняя компоненты --- также


кососимметрические тензоры, ортогональные между собой.


\end{props}





\begin{pf}


Разложение (12) непосредственно следует из представления тождественного


(1,1)-тензора по формуле (13)


\begin{gather}


\id =\pr +\ort,\\


\pr \circ \ort =\ort \circ \pr =0.


\end{gather}


Действительно, $Q=\id \circ Q\circ \id =(\pr +\ort )Q(\pr +\ort )=


(\pr Q\pr +\ort Q\ort )+(\pr Q\ort +\ort Q\pr )=|Q|+\{ Q\}$.


Если $Q$ кососимметричен, то воспользовавшись самосопряженностью


ортопроекторов $\pr $ и $\ort $, получим


$g(|Q|X,Y)=g(\pr Q\pr X+\ort Q\ort X,Y)=g(Q\pr X,\pr Y)+g(Q\ort X,\ort Y)=


-g(\pr X,$ $Q\pr Y)-g(\ort X,Q\ort Y)=-g(X,\pr Q\pr Y)-g(X,\ort Q\ort X)=


-g(X,|Q|Y)$.


Аналогично доказывается кососимметричность внешней компоненты. Ортогональность


внутренней и внешней компонент следует из предложения 1.2~а), свойства


идемпотентности ортопроекторов $\pr $ и $\ort $, а также равенств (14).


Действительно, $g(|Q|,\{ S\} )=-\Tr|Q|\{ S\}=0$.


\end{pf}





В дальнейшем ортонормальный репер $E$ будем называть {\it адаптированным} к


подпространству $\Pi$, если линейная оболочка первых $l$ векторов $E$ совпадает


с подпространством $\Pi$.





\begin{cors} Для матриц внутренней и внешней компонент


кососимметрических тензоров относительно репера $E$, адаптированного к


подпространству $\Pi$, имеют место представления


$$


|Q|=


\begin{pmatrix}


Q_l & 0 \\ 0 & Q_p


\end{pmatrix}


,\qquad


\{ Q\} =


\begin{pmatrix}


0 & -Q^t_\triangle\\


Q_\triangle & 0


\end{pmatrix}


,


$$


где подматрицы $Q_l$, $Q_p$, $Q_\triangle$ имеют размеры $l\times l$,


$p\times p$ и $l\times p$ соответственно.


\end{cors}





\begin{pf} Матрицы операторов проектирования на подпространство $\Pi$,


ортогональное дополнение к $\Pi$ и матрица кососимметрического


оператора $Q$ в данном


базисе $E$ будут иметь соответственно следующий вид:


$$


\pr =


\begin{pmatrix}


I_l & 0\\


0 & 0


\end{pmatrix}


,\qquad


\ort =


\begin{pmatrix}


0 & 0\\


0 & I_p


\end{pmatrix}


,\qquad


Q=


\begin{pmatrix}


Q_l & -Q^t_\triangle\\


Q_\triangle & Q_p


\end{pmatrix}


.


$$


Воспользовавшись определением 1.1 и правилом поблочного умножения матриц,


получим необходимый вид матриц внешней и внутренней компонент тензора $Q$.


\end{pf}





Приведенные ниже утверждения непосредственно вытекают из следствия 1.3.





\begin{cors} Для любых векторов $X$, $Y$ из касательного пространства


в точке $w$ многообразия $M$


\begin{align}


|X\land Y|&= \pr X\land \pr Y+\ort X\land \ort Y,\\


\{ X\land Y\} &=\ort X\land \pr Y+\pr X\land \ort Y.


\end{align}


\end{cors}





\begin{cors} Все пространство кососимметрических тензоров


$\boldsymbol\Lambda ^2_w M$ в точке $w\in M$ разлагается в прямую сумму


$\boldsymbol\Lambda^2_w M=I(\Pi)\oplus E(\Pi)$, где $\Pi$ --- некоторое


произвольное подпространство в $T_wM$.


\end{cors}





\begin{props} Для любых кососимметрических тензоров $Q$ и $S$


и любых векторов $X$, $Y$ и $Z$ на многообразии $M$ выполняются равенства


\begin{alignat*}{2}


&{\rm{а)}} \ |[Q,S]|=[|Q|,|S|]+[\{ Q\} ,\{ S\} ], &\qquad


&{\rm{в)}} \ \underset{(X,Y,Z)}{S}{}|X\land Y|(Z)=0,\\


&{\rm{б)}} \ \{ [Q,S]\} = [|Q|,\{ S\} ]+[\{ Q\},|S|], &\qquad


&{\rm{г)}} \ \underset{(X,Y,Z)}{S}{}\{ X\land Y\} (Z)=0.


\end{alignat*}


Все компоненты тензоров берутся относительно одного и того же подпространства


$\Pi\subset T_wM$, $w\in M$.


\end{props}





\begin{pf} а), б) Эти пункты являются прямым следствием правил


поблочного умножения матриц при вычислении коммутатора и следствия 1.3.


в) Из формулы (15) следует, что


\begin{gather*}


\underset{X,Y,Z}{S}{}|X\land Y|(Z)=\underset{X,Y,Z}{S}{}(\pr X\land


\pr Y)(Z)+\underset{X,Y,Z}{S}{}(\ort X\land \ort Y)(Z)=\\


=\underset{X,Y,Z}{S}{}(\pr X\land \pr Y)(\pr Z)+\underset{X,Y,Z}{S}{}


(\ort X \land \ort Y)(\ort Z).


\end{gather*}


Равенство нулю этого выражения следует из предложения 1.1~б),


примененного к векторам $\pr X$, $\pr Y$, $\pr Z$ и $\ort X$, $\ort Y$,


$\ort Z$.





г) Из предложения 1.4 и формулы (16) получим


$$


\underset{X,Y,Z}{S}{}\{ X\land Y\} (Z)=\underset{X,Y,Z}{S}{}X\land Y(Z)-


\underset{X,Y,Z}{S}{}|X\land Y|(Z).


$$


Равенство нулю этого выражения следует из предложения 1.1~б) и равенства


в) доказываемого предложения.


\end{pf}





\begin{cors} $[\{ Q\} ,\{ S\} ],[|Q|,|S|]\in I(\Pi )$,


$[\{ Q\} ,|S|],[|Q|,\{ S\} ]\in E(\Pi )$


$\forall Q,S\in \boldsymbol\Lambda ^2(M)$.


\end{cors}





\begin{props} Пусть $\Pi$ --- некоторое подпространство


касательного пространства к многообразию $M$ в точке $u$. Тогда операции


взятия внутренней и внешней компонент $(1,1)$-тензора относительно $\Pi$,


рассматриваемые как эндоморфизмы в пространстве кососимметрических


тензоров в точке $u$, являются ортопроекторами.


\end{props}





\begin{pf} Хорошо известно, что линейный оператор является


ортопроектором тогда, и только тогда, когда он является самосопряженным


идемпотентом. Таким образом, для доказательства предложения достаточно


показать, что выполняются равенства


\begin{align}


g(\{ Q\} ,S)&=g(Q,\{ S\} ),\\


\{ \{ Q\} \} &=\{ Q\} ,\\


g(|Q|,S)&=g(Q,|S|),\\


\|Q\|&=|Q|.


\end{align}


Из идемпотентности и самосопряженности операторов $\pr $ и $\ort $ в $T_uM$


следует, что


\begin{gather*}


g(\{ Q\} ,S)=g(\pr Q\ort ,S)+(\ort Q\pr ,S)=g(Q\ort ,\pr S)+g(Q\pr ,\ort S)=\\


=g(Q,\pr S\ort )+g(Q,\ort S\pr )=g(Q,\{ S\} ).


\end{gather*}


Равенство (19) доказывается аналогично, а


равенства (18), (20) вытекают из следствия 1.3.


\end{pf}





\section{Векторные поля в грассмановом расслоении}





Этот параграф посвящен построению образующей системы векторных полей на


грассмановом расслоении $G^l(M)$, относительно которой будет удобно установить


связь между геометрическими характеристиками базы $M$, слоя $G^l_n$ и


расслоения $G^l(M)$. При этом грассманово расслоение рассматривается как


риманово пространство с метрикой $\g g$, ассоциированной


посредством связности с метрикой базы $M$ и канонической симметрической


метрикой $\gamma $ грассманиана $G^l_n$ (\cite{Far1}, определение 2). Там же в


следствии 2 теоремы В для римановой метрики $\g g$ была


получена формула, которая с точностью до


введенного в \S\,1 понятия внешней компоненты имеет вид


\begin{equation}


\g g(Q_1,Q_2)=\pi ^\ast g(S_1,S_2)-\1 \Tr \{ \omega (S_1)


\} _{V_1}\circ \{ \omega (S_2)\} _{V_1}.


\end{equation}


Здесь $S_t$ ($t=1,2$) --- касательные векторы в расслоении ортонормальных


реперов


$O(M)$, прикрепленные в одной и той же произвольной точке, проектирующиеся на


касательные векторы $Q_t$ в грассмановом расслоении $G^l(M)$ c помощью проекции


$R_H$ ($O(M)$ --- можно рассматривать как главное расслоение над $G^l(M)$ со


структурной группой $H=O_l\times O_p$ (\cite{Kob1}, гл.\,1, предл.\,5.5)),


сопоставляющей


реперу из $n$ касательных векторов к $M$ $l$-мерное подпространство,


натянутое на первые $l$ векторов репера: $R_{H*}S_t=Q_t$; $\omega $


--- форма римановой связности на $O(M)$, $\pi $ --- проекция в расслоении


$O(M)$,


$V_1=L(e_1,\dotsc,e_l)$ --- линейная оболочка первых $l$ векторов канонического


базиса в $\Bbb R^n$.





Воспользуемся локальными координатами в $O(M)$ и $G^l(M)$, подробно


обсуждаемыми в \cite{Far1}, \cite{Far2}. В дальнейшем спектр изменения индексов $a, b, c, f$ от 1


до $n$; $i$, $j$, $k$, $m$ --- от 1 до $l$; $\alpha$, $\beta$ --- от 1 до


$p$ ($l$ и


$p=n-l$ фиксированные, $n=\dim M$). Матрицы $X$ размером $n\times n$


\begin{equation}


X=


\begin{pmatrix}


X_l & X^\triangle \\


X_\triangle & X_p


\end{pmatrix}


,\qquad


h=


\begin{pmatrix}


X^{-1}_l & 0\\


0 & X^{-1}_p


\end{pmatrix}


.


\end{equation}


где $X_l$, $X_p$ --- подматрицы размеров $l\times l$ и $p\times p$


соответственно, совместно с набором $w=(w^1,\dotsc,w^n)$


координат базового многообразия $M$ задают координаты точки в $O(M)$,


т.\,е. репера $EX=(E_aX^a_1,\dotsc,E_aX^a_n)$ ($E$ --- фиксированное


ортонормальное


реперное поле, определенное, вообще говоря, локально), а матрицы $\xi $ размером


$l\times p$ совместно с набором $w$ задают координаты точки в грассмановом


расслоении $G^l (M)$, т.\,е. координаты $l$-мерного подпространства


касательного пространства, являющегося решением системы уравнений из


1-форм в точке $w$:


\begin{equation}


\chi ^{l+\alpha }=\xi ^\alpha _k \chi ^k,


\end{equation}


где $\chi ^a$ --- кореперное поле, двойственное $E$. Естественно, последние


координаты определены лишь в окрестности подпространства $\Pi_0$, являющегося


линейной оболочкой $L(E_1,\dotsc,E_l)$ первых $l$ векторов репера $E$ в некоторой


точке $w_0$ базы $M$. Описанные координаты в $O(M)$ и $G^l(M)$ называются


\cite{Far1}, \cite{Far2} координатами относительно реперного поля $E$ (или кореперного поля


$\chi $). В последующем разбивка на блоки квадратных матриц размера


$n\times n$ принимается по умолчанию, как описано в (22).





\begin{lems} Проекция $R_H: O(M)\to G^l(M)$ в координатах


относительно реперного поля $E$ имеет следующее представление:


\begin{equation}


\xi = R_H(X)=X_\triangle \cdot X_l^{-1}


\end{equation}


\end{lems}





\begin{pf} Матрица $h$ из (22), которая определена в достаточно малой


окрестности $\Pi_0$, принадлежит структурной группе $H$.


Тогда, умножив матрицу $X$ из (22) на матрицу $h$, получим


$$


X\cdot h=


\begin{pmatrix}


X_l & X^\triangle \\


X_\triangle & X_p


\end{pmatrix}


\cdot


\begin{pmatrix}


X_l^{-1} & 0\\


0 & X_p^{-1}


\end{pmatrix}


=


\begin{pmatrix}


I_l & X^\triangle \cdot X_p^{-1}\\


X_\triangle \cdot X^{-1}_l & I_p


\end{pmatrix}


.


$$


Проверкой, аналогичной той, которая приведена в \cite{Far2}, нетрудно убедиться, что


репер с такими координатами обладает тем свойством, что


первые его $l$ векторов удовлетворяют системе (23).


\end{pf}





\begin{defns} Пусть $Q$ --- кососимметрическое тензорное


поле на $M$. Векторное поле $\olv Q$ на $O(M)$ назовем вертикальным лифтом


$Q$ в $O(M)$, если оно в каждой точке $u$ из $O(M)$ удовлетворяет системе уравнений


\begin{equation}


\pi _*\olv Q=0,\qquad \omega _u(\olv Q)=\ad u^{-1} Q=u^{-1}\circ Q\circ u.


\end{equation}


Здесь репер $u$ стандартным образом отождествляется с линейным оператором из


$\Bbb R^n$ на $T_{\pi (u)}M$.


\end{defns}





Приведенное определение 2.1 является частным случаем следующей общей


конструкции. Пусть $P(M,G)$ --- главное расслоение, $G$ --- его структурная


группа, а


$\frak g$ --- алгебра Ли группы $G$. Для эффективного левого действия $G$ на


$\frak g$ дифференциалами $\ad a_*$ построим ассоциированное расслоение $E$ со


стандартным слоем $\frak g$. Тогда для каждого сечения $Q$ расслоения $E$


уравнениями (26) однозначно определяется вертикальное векторное


поле $\olv Q$ в $P$


\begin{equation}


\pi _*\olv Q=0,\qquad u\left( \omega _u(\olv Q)\right) =Q.


\end{equation}


Здесь $\forall u\in P$ и $\forall A\in \frak g$ по определению ассоциированного


расслоения (\cite{Kob1}, с.\,60) $u(A)=\{ (ua,\ad a^{-1}_*A)\mid a\in G\} $


--- орбита


пары $(u,A)$, т.\,е. элемент из слоя $E_{\pi (u)}$. В случае, когда $G$


является


общей линейной группой $GL_n$ или ее подгруппой, как в определении 2.1,


элементы


из $E_x$ могут быть отождествлены с эндоморфизмами $T_xM$, т.\,е.


(1,1)-тензорами,


а вторая формула в (26) при этом превращается во вторую формулу в


(25), т.\,к. в этом случае $u(A)=u\circ A\circ u^{-1}$.





Следует отметить, что в любом главном расслоении имеется образующая система


вертикального распределения, состоящая из фундаментальных векторных полей


$A^*$,


(\cite{Kob1}, с.\,57), где как и прежде $A\in \frak g$. Однако от построенных


выше фундаментальное векторное поле отличается тем, что не является


правоинвариантным. А именно, правоинвариантность $\olv Q$, устанавливаемая ниже


в предложении 2.2, позволит корректно определить вертикальные векторные поля на


$G^l(M)$.


Традиционным образом определенное векторное поле $\olh Y$, являющееся


горизонтальным лифтом в $O(M)$ векторного поля $Y$ на $M$, удовлетворяет


системе уравнений


\begin{equation}


\pi _*\olh Y=Y,\qquad \omega (\olh Y)=0.


\end{equation}


Следующее утверждение дает эффективный способ конструирования вертикальных


лифтов из кососимметричеких тензоров.





\begin{props} В координатах $(w,X)$ на $O(M)$ относительно


ортонормального реперного поля $E$ векторное поле $\olv Q$ допускает


следующее представление:


\begin{equation}


\olv Q =\Tr QX\frac \partial{\partial X^t},


\end{equation}


где $Q$ -- матрица кососимметрического тензора относительно реперного


поля $E$, а $\frac \partial {\partial X^t}=(\frac \partial {\partial X})^t=


(\frac \partial {\partial X^a_b})$ --- матрица координатных вертикальных векторов.


\end{props}





\begin{pf} Пусть в точке $X$ из $O(M)$ вектор $\olv Q=F^a_b


\frac \partial {\partial X^a_b}=\Tr F\frac \partial {\partial X^t}$. Из


определения 2.1 следует, что $\omega _X(\olv Q)=X^{-1}QX$. Теперь из локально


координатного


представления формы связности $\omega = X^{-1}\delta X=X^{-1}(dX+Г(\chi )X)$


(где $Г(\chi )$ -- матрица Кристоффеля относительно кореперного поля $\chi $,


(\cite{Far1}, \S\,2)) в координатах $(w,X)$, находим матрицу $F=QX$.


\end{pf}





Непосредственной проверкой путем дифференцирования легко убедиться в


справедливости следующего следствия.





\begin{cors} Локальная 1-параметрическая группа локальных


преобразований, индуцируемая векторным полем $\olv Q$ на $O(M)$, имеет


координатное представление


$q_t(X)=\exp (tQX).$


\end{cors}





Очевидно, что в координатах $(w,X)$ относительно ортонормального реперного


поля $E$ горизонтальный лифт допускает следующее локально координатное


представление (\cite{Kob1}, c.\,140):


\begin{equation}


\olh Y=Y^a\frac \partial {\partial w^a}-Tr Г(Y)X\frac \partial


{\partial X^t},


\end{equation}


где $Г(Y)$ --- такая же матрица, как в (5), $Y^a$ --- компоненты векторного


поля $Y$ относительно координат базы $M$, а в матрице $\frac \partial


{\partial X^t}$ компонентами служат векторные поля, причем в $c$-й строке и


$a$-м столбце стоит вектор $\frac \partial {\partial X^a_c}$. Формула (29)


допускает представление, не использующее базовых координат, а именно,


$$


\olh Y=\chi^a(Y)[E_a]-\Tr Г(Y)X\frac \partial {\partial X^t},


$$


где $[E_a]$ --- координатное представление векторного поля $E_a$ в


произвольных координатах на $M$.





Риманова метрика $\ol g$ в $O(M)$, ассоциированная с римановой


метрикой базы $M$ и канонической симметрической метрикой стандартного


слоя $O_n$ (метрика Мока), может быть вычислена по формуле


\begin{equation}


\ol g(S_1,S_2)=\pi ^* g(S_1,S_2)-\1 \Tr \omega (S_1)\omega (S_2).


\end{equation}





\begin{lems} Пусть $Q$, $S$ --- кососимметрические тензорные поля, а


$X$, $Y$ --- векторные поля на римановом многообразии $M$. Тогда


\begin{gather}


\ol g(\olh X,\olh Y)=g(X,Y),\\


\ol g(\olv Q,\olh X)=0,\\


\ol g(\olv Q,\olv S)= -\1 \Tr QS =\1 g(Q,S),


\end{gather}


причем эти выражения не зависят от вертикальных координат, а след в формуле


{\em (33)} вычисляется в касательном пространстве к $M$.


\end{lems}





\begin{pf} Первые две формулы, очевидно, следуют из определения 2.1


и формулы (30). Докажем формулу (33).


\begin{gather*}


\ol g(\olv Q,\olv S)=-\1 \Tr \w (\olv Q)\w (\olv S)=-\1 \Tr


(u^{-1}\circ Q\circ u)(u^{-1}\circ S\circ u)=\\


=-\1 \Tr u^{-1}QSu=-\1 \Tr QS.


\end{gather*}


Последнее равенство записано в силу


коммутативности следного функционала. Применив п.~а) предложения 1.2, получим


вторую часть доказываемой формулы.


\end{pf}





\begin{props} Вертикальный лифт $\olv Q$ кососимметрического


тензорного поля $Q$ является правоинвариантным векторным полем на $O(M)$,


т.\,е. $R_{a*}\olv Q=\olv Q$ $\forall a\in O_n$. 


\end{props}





\begin{pf} Для любого элемента $a\in O_n$


$$


\w _{ua}(R_{a*}\olv Q)=(R_a^*\w )_u(\olv Q)=(\ad a^{-1})\w _u(\olv Q)=


(\ad a^{-1})(\ad u^{-1}Q)=


\ad (ua)^{-1}Q=\w _{ua}(\olv Q).


$$


Так как вертикальный лифт является вертикальным векторным полем и форма


связности $\w $ невырождена на вертикальных подпространствах, то это


равенство доказывает рассматриваемое предложение.


\end{pf}





Хорошо известно (\cite{Kob1}, гл.\,2, пр.\,1.2), что горизонтальный лифт


$\olh Y$ базового векторного


поля $Y$ есть правоинвариантное векторное поле на $O(M)$. Таким образом,


на расслоении ортонормальных реперов $O(M)$ построена (локально)


система правоинвариантных образующих в модуле векторных полей, т.\,е.


любое векторное поле на $O(M)$ может быть разложено в линейную


комбинацию $\1 n(n-1)$ вертикальных и $n$ горизонтальных лифтов.





\begin{lems} Для любых кососимметрических тензоров $Q$ и $S$


и любого векторного поля $Y$ на многообразии $M$ выполняются


равенства


\begin{align}


\olh Y\w (\olv Q)&=\w (\olv {(\n YQ)}),\\


\olv S\w (\olv Q)&=\w (\olv {[Q,S]}).


\end{align}


\end{lems}





{\bf Доказательство.} Пусть в $O(M)$ введены локальные координаты


$(w,X)$ относительно ортонормального реперного поля $E$. В ходе


доказательства предложения 2.1 было показано, что в локальных


координатах матрица функций $\w (\olv Q)$ на $O(M)$ имеет вид $\w


(\olv Q)=X^{-1}QX$, где $X$ --- матричные координаты точки, в которой


вычисляется значение формы связности. Применяя координатное выражение


(29) для горизонтального лифта и тот факт, что значение вертикального


вектора $\Tr C\ch X$ из $O(M)$ на координатной матрице $X$ равно $C$,


а на матрице $Q$ функций базовых координат равно 0, получим


координатное представление действия вектора на функцию в точке $(w,X)$:


\begin{multline*}


\olh Y\w (\olv Q)=(Y^a\chp{w^a}-\Tr Г(Y)X\ch X)(X^{-1}QX)=\\


=Y^a(X^{-1}\frac{\partial Q}{\partial w^a}X)-\Tr Г(Y)X\ch X(X^{-1}QX)=\\


=X^{-1}YQX-\Tr Г(Y)X\ch X(X^{-1})QX-X^{-1}Q\Tr Г(Y)X\ch X(X)=\\


=X^{-1}dQ(Y)X+X^{-1}\Tr Г(Y)X\ch X(X)X^{-1}QX-X^{-1}QГ(Y)X=\\


=X^{-1}\left( dQ(Y)+Г(Y)Q-QГ(Y)\right) X.


\end{multline*}


После применения формулы (5) и определения 2.1 получим требуемое равенство.


Из предложения 2.1 аналогично предыдущему получим


\begin{gather*}


\olv S\w (\olv Q)=\Tr SX\ch X(X^{-1}QX)=


(\Tr SX\ch X(X^{-1}))QX+X^{-1}Q(\Tr SX\ch X(X))=\\


=-X^{-1}SXX^{-1}QX+X^{-1}QSX=X^{-1}(QS-SQ)X=


X^{-1}[Q,S]X=\w (\olv {[Q,S]}).\quad\square


\end{gather*}





\begin{lems} Пусть $Q$, $S$ --- кососимметрические тензорные поля, а


$X$, $Y$ --- векторные поля на базе $M$. Тогда


\begin{align}


&{\rm{а)}} \ [\olv Q,\olv S] = -\olv {[Q,S]},\\


&{\rm{б)}} \ [\olh X,\olv Q] = \olv {(\n XQ)},\\


&{\rm{в)}} \ [\olh X,\olh Y] = \olh {[X,Y]} - \olv {R(X,Y)}.


\end{align}


\end{lems}





\begin{pf} а) Скобка Ли вертикальных векторных полей --- это


вертикальное векторное поле. Кроме того, из (\cite{Kob1}, c.\,43)


\begin{equation}


\w ([\olv Q,\olv S])=-2d\w (\olv Q,\olv S)+\olv Q\w (\olv S)-\olv S\w (\olv Q).


\end{equation}


Из структурного уравнения $d\w =-\1 [\w ,\w ]+\Omega$ (\cite{Kob1}, c.\,81) и


горизонтальности формы кривизны $\Omega$, т.\,е. равенства ее нулю при наличии


хотя бы одного вертикального вектора на основании леммы 2.3 следует, что


\begin{gather*}


\w ([\olv Q,\olv S])=[\w (\olv Q),\w (\olv S)]-2\Omega (\olv Q,\olv S)+\olv Q


\w (\olv S)-\olv S\w (\olv Q)=\\


=\w (\olv Q)\w (\olv S)-\w (\olv S)\w (\olv Q)+\olv Q\w (\olv S)-\olv S


\w (\olv Q)=\ad u^{-1}Q\circ \ad u^{-1}S-\\


-\ad u^{-1}S\circ \ad u^{-1}Q-\w (\olv {[Q,S]})-\w (\olv {[Q,S]})=


\ad u^{-1}[Q,S]-\\


-2\w (\olv {[Q,S]})=-\w (\olv {[Q,S]}).


\end{gather*}


Так как форма связности $\w$ невырождена на вертикальных векторах, то отсюда


следует равенство (36).





б) Отсутствие горизонтальной компоненты для $[\olh X,\olv Q]$ следует из того,


что векторные поля $\olh X$ и $\olv Q$ являются $\pi$-связанными с


векторными полями $X$ и нулевым соответственно. Из формулы (39), структурного


уравнения,


горизонтальности формы кривизны и вертикальности формы связности следует, что


$$


\w ([\olh X,\olv Q])=-2d\w (\olh X,\olv Q)+\olh X\w (\olv Q)-\olv Q\w(\olh X)=


\olh X\w (\olv Q)=\w (\olv {(\n XQ)}),


$$


причем последнее равенство записано на основании леммы 2.3.





в) Равенство $\pi_*([\olh X,\olh Y])=[X,Y]$ следует из


$\pi$-связности векторных полей $\olh X$, $\olh Y$ и $X$, $Y$


соответственно. Равенство $\w ([\olh X,\olh Y])=-2\Omega (\olh


X,\olh Y)$ и определение преобразования кривизны,


являющегося при фиксированных $X$ и $Y$ кососимметрическим


(1,1)-тензором, приводят к цепочке равенств: $\w ([\olh X, \olh


Y])=-u^{-1}R(X,Y)u=-\w (\olv {R(X,Y)})$. Отсюда следует формула (38).


\end{pf}





Введем понятие горизонтального лифта $\grh Y$ векторного поля


$Y$ на $M$ в грассманово расслоение $G^l(M)$ стандартной процедурой,


т.\,е. как горизонтальное векторное поле, проектирующееся в поле $Y$.


Это также можно сделать с помощью корректно определенного равенства


\begin{equation}


R_{H*}\olh Y=\grh Y,


\end{equation}


т.\,к. $R_H$ --- отображение связности (\cite{Far1}, из доказательства


теоремы В), а поле $\olh Y$,


как отмечалось выше, правоинвариантное. Таким образом, горизонтальные


лифты в $O(M)$ и $G^l(M)$ одного и того же векторного поля на


базе $M$, $R_H$-связны между собой.





\begin{defns} Пусть $Q$ --- кососимметрическое


тензорное поле на $M$. Векторное поле $\grv Q$


на $G^l(M)$ назовем вертикальным лифтом $Q$ в грассманово


расслоение, если оно получено проекцией посредством $R_H$ из


вертикального лифта $\olv Q$ того же тензорного поля $Q$ в


$O(M)$, т.\,е. $\grv Q=R_{H*}\olv Q$.


\end{defns}





Из предложения 2.2. следует, что вертикальный лифт $\grv Q$ в


$G^l(M)$ существует для любого кососимметрического тензора $Q$.


Очевидно, что в окрестности любой точки из $G^l(M)$ можно построить


систему векторных полей из горизонтальных и вертикальных лифтов,


порождающую (локально) весь модуль векторных полей в $G^l(M)$, взяв, например,


$\grh {(\chp {w^a})}$ и $\grv {(\chp {w^a}\land \chp {w^b})}$. Необходимо


отметить, что из предъявленной системы образующих полей $n$


горизонтальных лифтов являются линейно независимыми в каждой точке и


порождают горизонтальное распределение в $G^l(M)$, а среди вертикальных


лифтов в точности $lp$ векторов линейно независимы в каждой точке


грассманова расслоения.





\begin{lems} Если кососимметрическая матрица $Q$ имеет


блочно-диагональный вид


\begin{equation}


Q=


\begin{pmatrix}


Q_l & 0 \\


0 & Q_p


\end{pmatrix}


,


\end{equation}


то ортогональная матрица ее экспоненты $e^Q$ имеет


также блочно-диагональный вид


\begin{equation}


e^Q =


\begin{pmatrix}


e^{Q_l} & 0\\


0 & e^{Q_p}


\end{pmatrix}


.


\end{equation}


\end{lems}





\begin{pf} Так как матрица $Q=Q_1+Q_2$, где


$$


Q_1=


\begin{pmatrix}


Q_l & 0 \\


0 & 0


\end{pmatrix}


,\qquad


Q_2=


\begin{pmatrix}


0 & 0\\


0 & Q_p


\end{pmatrix}


,


$$


а слагаемые обладают тем свойством, что $Q_1 Q_2=Q_2 Q_1=0$, то из


определения экспоненты следует, что $e^Q=e^{Q_1+Q_2}=e^{Q_1}e^{Q_2}$.


\end{pf}





\begin{props} Пусть $R_Hu=\Pi $. Тогда вертикальным


подпространством римановой субмерсии $R_H$ в точке $u$ служит


вертикальный лифт внутреннего относительно $\Pi$ подпространства


кососимметрических тензоров, т.\,е. $\Ker R_{H*}=\olv {I(\Pi )}$, где


$\olv {I(\Pi )}=\{ \olv {Q_u}, \{ Q\} _\Pi =0\} $, а горизонтальным


подпространством


римановой субмерсии $R_H$ в точке $u$ служит $H_uO(M)+\olv {E(\Pi )}$


прямая ортогональная сумма горизонтального {\em (}в смысле


римановой связности{\em )}


пространства $H_uO(M)$ и вертикального лифта внешнего относительно $\Pi $


подпространства кососимметрических тензоров $\olv {E(\Pi )}=\{ \olv {Q_u},


|Q|_\Pi =0\} $.


\end{props}





\begin{pf} Из следствия 2.1 и леммы 2.5 вытекает,


что все векторы из $\olv {I(\Pi )}$ касаются слоя через $u$. Кроме


того, размерность подпространства $I(\Pi )$, равная $\1 (l(l-1)+p(p-1))$,


совпадает с размерностью группы $H$. Значит, $\olv


{I(\Pi )}$ действительно является вертикальным подпространством римановой


субмерсии $R_H$. Касательное пространство в точке $u$ к риманову


многообразию $O(M)$ с метрикой $\ol g$, ассоциированной с метриками слоя и


базы, является прямой ортогональной суммой $H_uO(M)+V_uO(M)$


горизонтального и вертикального (в смысле римановой связности)


подпространств. В свою очередь $V_uO(M)$ является прямой суммой $\olv


{I(\Pi )}+\olv {E(\Pi )}$, что следует непосредственно из предложения 1.4 и


определения 2.1. Ортогональность $\olv {I(\Pi )}$ и $\olv {E(\Pi )}$ следует из


леммы 2.2:


$$


\ol g(\olv {\{ Q\} },\olv {|Q|})=\1 g(\{ Q\} ,|Q|)=0.


$$


Все это и доказывает данное предложение.


\end{pf}





\begin{cors}


$\grv {Q_\Pi }=\grv {(\{ Q\} _\Pi )}$.


\end{cors}





Итак, задача построения системы образующих векторных полей в $G^l(M)$


решена. Следующее предложение даст эффективный способ конструирования


таких лифтов, после чего можно будет перейти к решению основной задачи:


установлению связи между геометрией грассманова расслоения и базы.





\begin{props} В локальных координатах $(w,\xi )$ в $G^l(M)$


относительно кореперного поля $\chi $ вертикальный лифт тензорного поля


$Q$ имеет следующее представление:


\begin{equation}


\grv Q=\Tr (Q_\tri +Q_p\xi -\xi Q_l-\xi Q^\tri \xi )\ch \xi,


\end{equation}


где матрица компонент $Q$ относительно $\chi $ имеет блочную структуру


обозначений, как в {\em (22)}.


\end{props}





\begin{pf} Продифференцируем формулу (24) леммы 2.1:


\begin{multline*}


d\xi = dX_\tri X_l^{-1}+X_\tri dX_l^{-1}=dX_\tri X_l^{-1}-X_\tri


X_l^{-1}dX_lX_l^{-1}=\\


=(dX_\tri -X_\tri X_l^{-1}dX_l)X_l^{-1}=(dX_\tri -\xi dX_l)X_l^{-1}.


\end{multline*}


Так как из поблочного умножения матриц следует, что


$$


QX=


\begin{pmatrix}


Q_l & Q^\tri \\


Q_\tri & Q_p


\end{pmatrix}


\cdot


\begin{pmatrix}


X_l & X^\tri \\


X_\tri & X_p


\end{pmatrix}


=


\begin{pmatrix}


(QX)_l & (QX)^\tri \\


(QX)_\tri & (QX)_p


\end{pmatrix}


,


$$


где $(QX)_l=Q_lX_l+Q^\tri X_\tri $, $(QX)_\tri =Q_\tri X_l+Q_pX_\tri $,


то из предложения 2.1 вытекает, что $dX_l(\olv Q)=(QX)_l$, $dX_\tri (\olv


Q)=(QX)_\tri $.


Таким образом,


\begin{gather*}


d\xi(\olv Q)=(dX_\tri -\xi dX_l)X_l^{-1}(\olv Q)=((QX)_\tri -


\xi (QX)_l)X_l^{-1}=(Q_\tri X_l+\\


+Q_pX_\tri -\xi Q_lX_l-\xi Q^\tri X_\tri)X_l^{-1}=Q_\tri +Q_pX_\tri


X_l^{-1}-\xi Q_l-\xi Q^\tri X_\tri X_l^{-1}=\\


=Q_\tri +Q_p\xi -\xi Q_l-\xi Q^\tri \xi.


\end{gather*}


Значит, $\grv Q=R_{H*}(\olv Q)=\Tr (Q_\tri +Q_p\xi -\xi Q_l-\xi Q^\tri


\xi )\ch {\xi}$.


\end{pf}





Доказанное предложение позволяет видеть,что в точке $\Pi $ с координатами


$\xi =0$ вертикальный лифт кососимметрического тензора $Q$ зависит


лишь от $lp$ компонент углового блока $Q_\tri $ матрицы $Q$.


Итак, задача построения системы образующих векторных полей в $G^l(M)$


решена.





\begin{lems} Пусть $Q$ и $S$ --- кососимметрические тензорные поля, а


$X$ и $Y$ --- векторные поля на $M$. Тогда справедливы следующие тождества:


\begin{align*}


&{\rm {а)}} \ [\grv Q,\grv S]=-\grv {[Q,S]},\\


&{\rm {б)}} \ [\grh X,\grv Q ]=\grv{(\n XQ)},\\


&{\rm {в)}} \ [\grh X,\grh Y]=\grh {[X,Y]}-\grv {R(X,Y)}.


\end{align*}


\end{lems}





\begin{pf} Как было показано выше, векторные поля $\olv Q$


и $\grv Q$, а также $\olh X$ и $\grh X$ попарно $R_H$-связны. Тогда


выполнение тождеств леммы 2.4 влечет за собой выполнение соответствующих


тождеств этой леммы.


\end{pf}





\begin{lems} Пусть $Q$, $S$ --- кососимметрические тензорные


поля, а $X$, $Y$ --- векторные поля на $M$. Тогда для их


вертикальных и горизонтальных лифтов в $G^l(M)$ верны следующие


тождества:


\begin{align}


&\g g(\grh X,\grh Y)=g(X,Y),\\


&\g g(\grv Q,\grh X)=0,\\


&\g g_\Pi (\grv Q,\grv S)= -\1 \Tr \{ Q\} \{ S\} =\1 g(\{ Q\} , \{ S\}),


\end{align}


где внешние компоненты тензорных полей берутся относительно


подпространства в $G^l(M)$, в котором вычисляется $\g g$.


\end{lems}





{\bf Доказательство.} Так как $R_H$ -- риманова субмерсия $O(M)$


на $G^l(M)$, а векторные поля $\olv Q$ и $\grv Q$, $\olh X$ и $\grh X$


попарно $R_H$-связны, то из предложения 2.3 и аналогичных доказываемым


тождеств леммы 2.2 немедленно следуют формулы (44) и (45). Для


доказательства тождества (46) обратимся к формуле (21) для вычисления


метрики $\g g$ в $G^l(M)$ и определениям 2.1 и 2.2. Тогда


$$


\g g_\Pi (\grv Q,\grv S)=g(\pi _*\olv Q,\pi _*\olv S)-\1 \Tr\{ \w (\olv


Q)\}_{V_1}\{ \w (\olv S)\}_{V_1}=


-\1 \Tr \{ u^{-1}Qu\} _{V_1}\circ \{u^{-1}Su\} _{V_1}.


$$


% tts the formula ended by = sign, repl. by "."; cf. formula below


При этом репер $u$ из $O(M)$ связан с $l$-мерным подпространством


$\Pi \in G^l(M)$ соотношением $R_Hu=\Pi $. Кроме того, линейное


отображение $u$ переводит подпространство $V_1$ --- линейную оболочку


первых $l$ векторов канонического базиса $\Bbb R^n$ --- в подпространство $\Pi$


из $T_xM$ $(x=\pi (u))$, а его ортогональное дополнение


$V_2=L(e_{l+1},\dotsc,e_n)$ --- в ортогональное дополнение к $\Pi $ в $T_xM$.


Следовательно, $\{ \ad u^{-1}Q\} _{V_1}=u^{-1}\{ Q\} _\Pi u=\ad u^{-1}\{


Q\} _\Pi$. Теперь, воспользовавшись коммутативностью следного функционала и


предложением 1.2~а), продолжим равенство


$$


-\1 \Tr(u^{-1}\{ Q\}_\Pi u)(u^{-1}\{ S\}_\Pi u)=\1 g(\{ Q\}_\Pi,\{ S\}_\Pi).


\quad \square


$$





В отличиe от вертикальных лифтов в расслоении ортонормальных реперов, в


выражении (46) есть зависимость от вертикальных координат. Это будет


прояснено в доказательстве следующей леммы.





\begin{lems} Если $(w,X)$ --- координаты репера $u$ {\em (}относительно


реперного поля $E)$ проектирующегося посредством $R_H$ в


подпространство $\Pi $, а $Q$, $S$ --- некоторые кососимметрические


тензорные поля на $M$, то имеет место координатное представление


\begin{equation}


\g g_\Pi (\grv Q,\grv S)=-\Tr I_{p,n}X^tQXI_{l,n}X^tSX,


\end{equation}


где координатные матрицы $n\times n$ имеют следующий вид:


$$


I_{p,n}=


\begin{pmatrix}


0 & 0 \\


0 & I_p


\end{pmatrix}


,\qquad


I_{l,n}=


\begin{pmatrix}


I_l & 0\\


0 & 0


\end{pmatrix}


,


$$


а $I_p$, $I_l$ --- единичные матрицы порядков $p$ и $l$ соответственно.


\end{lems}





\begin{pf} Так как матрица проектирования в $\Bbb R^n$ на


подпространство $V_1$ в каноническом базисе --- это $I_{l,n}$, а на


подпространство $V_2$ --- $I_{p,n}$, то из координатного представления


матрицы $\w (\olv Q)$ и определения 1.1 получим


$$


\{ \w (\olv Q)\} _{V_1}=\pr \w (\olv Q)\ort +\ort \w (\olv Q)\pr =


I_{l,n}X^{-1}QXI_{p,n}+I_{p,n}X^{-1}QXI_{l,n}.


$$


Используя свойства ортопроекторов $I_{l,n}\cdot I_{l,n}=I_{l,n}$ и


$I_{p,n}\cdot I_{p,n}=I_{p,n}$, а также $I_{l,n}\cdot I_{p,n}=I_{p,n}\cdot


I_{l,n}=0$ и формулу (21), запишем


\begin{gather*}


\g g_\Pi (\grv Q,\grv S)=-\1 \Tr (I_{p,n}X^{-1}QXI_{l,n}+I_{l,n}X^{-1}


QXI_{p,n})(I_{p,n}X^{-1}SXI_{l,n}+\\


+I_{l,n}X^{-1}SXI_{p,n})=-\1 \Tr (I_{p,n}X^{-1}QXI_{l,n}X^{-1}SXI_{p,n}


+I_{l,n}X^{-1}SXI_{p,n}\\


X^{-1}QXI_{l,n})=-\1 \Tr(I_{p,n}X^{-1}QXI_{l,n}X^{-1}SX+I_{l,n}X^{-1}


SXI_{p,n}X^{-1}QX)=\\


=-\Tr I_{p,n}X^{-1}QXI_{l,n}X^{-1}SX=-\Tr I_{p,n}X^tQXI_{l,n}X^tSX.


\end{gather*}


Последнее равенство получено из ортогональности матрицы $X$.


\end{pf}





\begin{thebibliography}{9} 





\bibitem{Shap}


{Шапуков Б.Н.}


{\em Связности на дифференцируемых расслоениях} //


Итоги науки и техн. ВИНИТИ.


Пробл. геометрии. -- 1983. -- Т.\,15. -- С.\,61--93.





\bibitem{Bor}


{Борисенко А. А., Ямпольский А. Л.}


{\em Риманова геометрия расслоений} //


УМН. -- 1991. -- Т.\,46. -- Вып.\,6. -- С.\,51--95.





\bibitem{Kob1}


{Кобаяси Ш., Номидзу К.}


{\em Основы дифференциальной геометрии}. Т.\,1. -- 


M.: Наука, 1981, -- 344~c.





\bibitem{Kob2}


{Кобаяси Ш., Номидзу К.}


{\em Основы дифференциальной геометрии}. Т.\,2. -- 


M.: Наука, 1981, -- 414~c.





\bibitem{Far1}


{Фарафонова Н.К.}


{\em Римановы метрики в расслоенных пространствах.} I //


Изв. вузов. Математика. -- 1995. -- \No\,7. -- C.\,65--77.





\bibitem{Far2}


{Фарафонова Н. К.}


{\em Римановы метрики в расслоенных пространствах.} II.


{\em Метрика грассманова расслоения}//


Изв. вузов. Математика. -- 1996. -- \No\,2. -- С.\,59--72.





\end{thebibliography} 


 


\vskip 2cm 


 


\postup{\it Харьковский государственный университет}{$\qquad$\it Поступили}


\postup{}{{\it первый вариант} 03.07.1995} 


\postup{}{{\it окончательный вариант} 22.08.1996} 


 


\end{document} 


